%
% 第3章：损失函数

\clearpage
\cleardoublepage

\chapter{损失函数}

\section{引言}

损失函数（也称为目标函数、成本函数或误差函数）是量化神经网络预测与真实目标匹配程度的数学度量。训练过程旨在通过基于梯度的优化调整网络参数来最小化这个损失。

损失函数应该：
\begin{enumerate}
    \item 可微（用于基于梯度的优化）
    \item 适当地惩罚较大误差
    \item 鼓励特定任务的正确行为
    \item 具有有用的统计解释
\end{enumerate}

\section{回归损失函数}

回归任务预测连续值。

\subsection{均方误差（MSE）}

Mean Squared Error（L2损失）：
\begin{equation}
\mathcal{L}_{MSE} = \frac{1}{N}\sum_{i=1}^{N} \|\mathbf{y}_i - \hat{\mathbf{y}}_i\|^2
\end{equation}

属性：
\begin{itemize}
    \item \textbf{凸函数（线性模型）}
    \item \textbf{对异常值敏感}
    \item \textbf{假设高斯误差分布}
\end{itemize}

\begin{advantages}
    \item 凸函数，易于优化
    \item 闭式解（线性回归）
    \item 平滑梯度
\end{advantages}

\begin{disadvantages}
    \item 对异常值敏感
    \item 二次惩罚可能导致大误差主导
\end{disadvantages}

\subsection{平均绝对误差（MAE）}

Mean Absolute Error（L1损失）：
\begin{equation}
\mathcal{L}_{MAE} = \frac{1}{N}\sum_{i=1}^{N} \|\mathbf{y}_i - \hat{\mathbf{y}}_i\|_1
\end{equation}

对异常值鲁棒，但在零处不可微。

\section{分类损失函数}

分类任务预测离散类别标签。

\subsection{二元交叉熵}

对于二分类 ($y \in \{0, 1\}$)：
\begin{equation}
\mathcal{L}_{BCE} = -\frac{1}{N}\sum_{i=1}^{N}[y_i \log \hat{y}_i + (1-y_i) \log(1-\hat{y}_i)]
\end{equation}

\begin{advantages}
    \item 概率解释
    \item 等价于二分类的最大似然估计
    \item 即使预测自信但错误时也有强梯度
\end{advantages}

\begin{code}
    \captionof{listing}{PyTorch中的损失函数}
    \label{code:loss_chinese}
\begin{minted}{python}
import torch.nn as nn

# 分类
criterion = nn.CrossEntropyLoss()  # 多分类
criterion = nn.BCEWithLogitsLoss()  # 二分类

# 回归
criterion = nn.MSELoss()  # MSE
criterion = nn.L1Loss()   # MAE
criterion = nn.SmoothL1Loss()  # Huber loss
\end{minted}
\end{code}

\section{高级损失函数}

\subsection{Focal Loss}

解决密集目标检测中的类别不平衡：
\begin{equation}
\mathcal{L}_{\text{Focal}} = -\alpha_t (1 - \hat{y}_t)^\gamma \log(\hat{y}_t)
\end{equation}

其中 $\gamma \geq 0$ 聚焦于难以分类的样本。

\subsection{标签平滑}

通过软化硬标签进行正则化：
\begin{equation}
y_{LS} = y(1 - \epsilon) + \frac{\epsilon}{K}
\end{equation}

其中 $\epsilon \in [0, 1]$ 是平滑参数。

\begin{advantages}
    \item 正则化效果
    \item 防止过度自信的预测
    \item 改善校准
\end{advantages}

\section{选择损失函数}

\begin{table}[h]
    \centering
    \caption{损失函数选择指南}
    \begin{tabular}{L{2.5cm} L{4cm} L{4cm}}
        \toprule
        \textbf{任务} & \textbf{推荐损失} & \textbf{备注} \\
        \midrule
        二分类 & BCE & 标准选择 \\
        多分类 & 交叉熵 & 可选标签平滑 \\
        回归 & MSE & 对异常值敏感 \\
        鲁棒回归 & MAE/Huber & 抗异常值 \\
        类别不平衡分类 & Focal Loss & 聚焦困难样本 \\
        分割 & Dice + CE & 处理类别不平衡 \\
        \bottomrule
    \end{tabular}
\end{table}

\section{本章小结}

\begin{enumerate}
    \item \textbf{损失函数量化预测误差} 并指导网络优化。
    \item \textbf{MSE} 是回归的标准但对异常值敏感。
    \item \textbf{交叉熵} 是分类的标准，具有概率解释。
    \item \textbf{高级损失} 如Focal Loss和Dice Loss解决特定挑战。
    \item \textbf{标签平滑} 提供正则化和改善校准。
\end{enumerate}
